\section{Experiments}
\label{sec:experiments}

We evaluate SART from a mechanism-first perspective. Rather than focusing only on end accuracy, we ask four questions: \textbf{(Q1)} can SART improve robustness under controlled shortcut settings where the ground-truth reasoning rule is known; \textbf{(Q2)} is directional correction---rather than detection accuracy alone---the load-bearing component; \textbf{(Q3)} do the gradient signatures underlying \textsc{ShortcutScore} generalize beyond the synthetic regime; and \textbf{(Q4)} what are the limitations of applying directional correction to pretrained autoregressive LLMs. This organization separates what SART reliably achieves today from what requires further adaptation at scale.

\subsection{Experimental Setup}
\label{sec:exp_setup}

\textbf{Models and training.} We use a GPT-style transformer (4 layers, $d_{\mathrm{model}}=256$, 8 heads, $d_{\mathrm{ff}}=1024$, $\sim$3.2M params) trained for 40 epochs with AdamW~\cite{loshchilov2019decoupled} (lr $1\!\times\!10^{-3}$, weight decay $10^{-4}$, batch 64, cosine schedule). SART hyperparameters $(\alpha=\beta=1.0,\tau_A=0.3,\tau_R=0.5,\gamma=1.0,\rho=0.3)$ are picked by an 80-configuration grid (Appendix~\ref{sec:appendix_sensitivity}).

\textbf{Datasets.} Three synthetic reasoning benchmarks under controlled shortcut injection: 2{,}000 train / 500 val / 1{,}000 test, training 70\% shortcut-rule and 30\% true-rule, val/test follow the true rule, test split evenly into a \emph{clean} half (shortcut and true rule agree) and a \emph{perturbed} half (the shortcut feature contradicts the correct label). \textbf{Math-Arithmetic}: $a{+}b{\geq}10$; shortcut $a{\geq}5$. \textbf{Financial-Analysis}: regulatory compliance from four features; shortcut revenue${\geq}5$, true rule margin${\geq}5$ \emph{and} debt${<}5$. \textbf{Causal-Reasoning}: causal-direction inference with a confounder; shortcut correlation${\geq}5$, true rule $x{\geq}5$ \emph{and} $z{<}3$.

\textbf{Baselines.} Eleven methods in four families. \emph{Loss / inference-time}: \textbf{SFT}, \textbf{Self-Consistency}~\cite{wang2023selfconsistency}, \textbf{Focal Loss}~\cite{lin2018focallossdenseobject}, \textbf{JTT}~\cite{liu2021justtraintwiceimproving}. \emph{Data-centric}: \textbf{Data Filtering} (drop confidence$>0.90$ post-warmup), \textbf{Influence Filtering}~\cite{koh2020understandingblackboxpredictionsinfluence}. \emph{Invariance / worst-group}: \textbf{Group DRO}$^\dagger$~\cite{sagawa2020distributionallyrobustneuralnetworks}, \textbf{IRM}$^\ddagger$~\cite{arjovsky2020invariantriskminimization}, \textbf{V-REx}$^\ddagger$~\cite{krueger2021outofdistributiongeneralizationriskextrapolation}, \textbf{Fishr}$^\ddagger$~\cite{rame2022fishrinvariantgradientvariances}. \emph{Failure-mining}: \textbf{LfF}~\cite{nam2020learningfailuretrainingdebiased}. $^\dagger$ requires group annotations; $^\ddagger$ multiple environments.

\textbf{Metrics.} \textbf{Clean Accuracy} ($\uparrow$) on the clean split; \textbf{Robustness} ($\uparrow$) on the perturbed split, where a model that has truly learned the rule should be unaffected and a shortcut-reliant one flips; \textbf{Reasoning Consistency} ($\uparrow$) on intermediate-token match; \textbf{Detection F1} ($\uparrow$) over a 50-threshold sweep; and the gradient-alignment statistic $|\cos(\mathbf{g}_s,\mathbf{g}_{\mathcal{V}})|$ over up to 100 held-out samples.

\subsection{Controlled Shortcut Benchmarks (Q1)}
\label{sec:exp_main}

\begin{table}[t]
\centering
\caption{Main results across three synthetic benchmarks. \textbf{Bold}: best; \underline{underline}: second best. Top: averages across all three benchmarks. Bottom: per-dataset robustness for representative methods (per-dataset accuracy and reasoning consistency in Appendix~\ref{sec:appendix_per_dataset}). $^\dagger$Requires group annotations. $^\ddagger$Requires multiple environments.}
\label{tab:main_results}
\footnotesize
\setlength{\tabcolsep}{4pt}
\begin{tabular}{lcccccccccccc}
\toprule
\textbf{Metric} & SFT & \makecell{Self-\\Cons.} & \makecell{Data\\Filt.} & JTT & \makecell{Focal\\Loss} & \makecell{DRO\\$^\dagger$} & \makecell{IRM\\$^\ddagger$} & \makecell{V-REx\\$^\ddagger$} & \makecell{Fishr\\$^\ddagger$} & LfF & \makecell{Infl.\\Filt.} & \textbf{SART} \\
\midrule
\multicolumn{13}{l}{\emph{Averages across three benchmarks}} \\
Acc.\ ($\uparrow$)       & 68.8 & 69.1 & 74.7 & 67.9 & 68.2 & 73.5 & 73.3 & 72.1 & 68.1 & 69.0 & \underline{76.0} & \textbf{80.7} \\
Rob.\ ($\uparrow$)       & 19.4 & 23.3 & \underline{54.6} & 19.0 & 18.2 & 36.0 & 34.4 & 34.9 & 16.7 & 19.4 & 47.7 & \textbf{54.8} \\
Reasoning ($\uparrow$)   & 55.9 & 55.9 & 64.7 & 54.9 & 55.4 & 61.5 & 61.5 & 60.6 & 56.3 & 56.3 & \underline{68.4} & \textbf{75.4} \\
SC Det.\ F1 ($\uparrow$) & --   & --   & \underline{0.69} & --   & --   & --   & --   & --   & --   & --   & --   & \textbf{0.82} \\
\midrule
\multicolumn{13}{l}{\emph{Per-dataset robustness ($\uparrow$, representative methods)}} \\
Math Rob.        & 3.2  & 11.1 & 26.6 & 0.0  & 0.0  & \underline{48.0} & 45.7 & 48.7 & 0.0  & 0.0  & \textbf{50.4} & 41.9 \\
Financial Rob.   & 28.6 & 11.7 & \textbf{68.4} & 10.2 & 23.4 & 31.2 & 46.0 & 50.4 & 0.4  & 1.7  & 23.6 & \underline{43.0} \\
Causal Rob.      & 26.4 & 11.3 & 68.8 & 1.5  & 1.3  & 28.8 & 44.2 & 45.9 & 0.9  & 0.4  & \underline{69.2} & \textbf{79.5} \\
\bottomrule
\end{tabular}
\end{table}

\textbf{Main results.} Table~\ref{tab:main_results} shows SART consistently improves both accuracy and robustness across all three datasets, reaching 80.7\% / 54.8\% versus the strongest baseline Influence Filtering at 76.0\% / 47.7\%---a $+4.7$\,pp accuracy and $+7.1$\,pp robustness gain---and improving over plain SFT by $+11.9$\,pp accuracy and $+35.4$\,pp robustness, all without group annotations or multiple environments.

\textbf{Interpretation.} The improvement is not driven by better loss optimization but by modifying the direction of updates. Loss-shaping baselines see only scalar loss values and cannot tell shortcut samples apart from legitimately-hard ones; filtering baselines drop entire samples and destroy validation-aligned residuals; annotation-based baselines require side information SART does not consume. SART operates at the gradient level, where the harmful and useful components within a single sample become explicit.

\textbf{Where SART helps most.} The largest gains appear on Causal-Reasoning ($79.5$, $+10.3$\,pp over Influence Filtering), where shortcut and true rules produce similar loss values but different gradient directions. On single-feature Math, simple confidence-based filtering is already strong (Influence Filtering $50.4$ vs.\ SART $41.9$), leaving less room for directional correction; on conjunction-rule Financial, Data Filtering's confidence threshold catches the shortcut directly ($68.4$ vs.\ SART $43.0$). The pattern supports our central claim: \emph{SART provides the most benefit when shortcut and reasoning signals are indistinguishable at the loss level but separable in gradient geometry}.

\subsection{Ablation: What Component Matters? (Q2)}
\label{sec:exp_ablation}

We compare SART against three variants. \textbf{SART-w/o-projection} disables conflict-only projection and keeps only answer-token suppression. \textbf{SART-w/o-suppress} sets $\rho=0$ and keeps only conflict-only projection. \textbf{SART-rwt} replaces directional surgery with scalar exponential reweighting $w(s)=\exp(-\lambda S(s))$ at $\lambda=3.0$, holding the detection step fixed.

\textbf{Findings.} Removing conflict-only projection collapses robustness from $54.8$ to $18.3$\,pp ($-36.5$\,pp), while removing answer-token suppression has only minor effect ($-3.9$\,pp robustness, $-1.1$\,pp accuracy). Scalar reweighting reaches the highest standalone detection F1 of $0.85$ yet lags full SART by $-9.3$\,pp accuracy and $-14.0$\,pp robustness. The mean alignment $|\cos(\mathbf{g}_s,\mathbf{g}_{\mathcal{V}})|$ tracks the same ordering: SFT $0.08$, SART-rwt $0.06$, w/o-proj $0.07$, w/o-suppress $0.04$, SART (Full) $\mathbf{0.03}$.

\textbf{Interpretation.} Detection alone is insufficient; the key factor is how detection is acted upon. Scalar reweighting fails because magnitude-only correction never changes update direction---a downweighted shortcut gradient still points the wrong way. Directional correction is the only mechanism that removes harmful gradient components while preserving the useful signal flagged samples still carry.

\subsection{Geometry Diagnostics (Q3)}
\label{sec:exp_detection}

\begin{figure}[t]
    \centering
    \includegraphics[width=\linewidth]{results/figure3.png}
    \caption{Detection-then-correction evidence chain. \textbf{(a)} \textsc{ShortcutScore} $S(s)$ vs.\ ground-truth shortcut rate, with per-sample raster, 20-bin means, and least-squares fit; $S(s)$ ranks samples correctly. \textbf{(b)} Alignment $A(s)$ stratified into TP/FP/TN/FN cells against the F1-optimal threshold $\tau_S^\star$; the per-sample conflict gate filters within the rank. \textbf{(c)} Clean accuracy and robustness across SART variants; conflict-only projection is the single load-bearing primitive.}
    \label{fig:analysis}
\end{figure}

We analyze whether \textsc{ShortcutScore} captures meaningful structure in gradient space.

\textbf{Findings.} \textsc{ShortcutScore} reaches an F1 of $0.82$ (Fig.~\ref{fig:analysis}b), comparable to the best discrete-detector baseline, and ranks shortcut samples consistently above clean ones. Importantly, the false-positive cell in panel (b)---samples flagged as shortcut by $S(s)$ but ground-truth clean---shows mean $A(s)=+0.02$, with $\sim$39\% still carrying positively-aligned signal at $\tau_S^\star$.

\textbf{Interpretation.} Shortcut detection is inherently a ranking problem rather than a binary classification problem. The downstream conflict gate $\mathbf{g}_s^\top\mathbf{g}_{\mathcal{V}}<0$ plays a critical role by selectively applying correction only when gradients truly conflict with validation improvement, recovering positively-aligned signal that the detector would otherwise discard.

\textbf{Hyperparameter sensitivity.} The projection strength $\gamma$ should scale with shortcut complexity: $\gamma=0.8$ best for single-feature Math (full projection over-erases useful residual), $\gamma=1.0$ best for multi-feature Financial and Causal. The suppression coefficient $\rho$ is inactive unless $R(s)>\tau_R$. Default $(\gamma=1.0,\rho=0.3)$ degrades gracefully; per-task tuning adds $+8.2$\,pp robustness on average (Appendix~\ref{sec:appendix_sensitivity}).

\subsection{Transfer to Pretrained LLMs (Q4)}
\label{sec:exp_llm_transfer}

We test SART on Qwen3-8B with LoRA rank-16 adapters on $q/k/v/o$ projections, using a controlled-injection pilot on GSM8K (70\% sum-of-numbers shortcut relabeling, evaluated on original-answer splits).

\textbf{Detection transfers.} \textsc{ShortcutScore} separates injected samples (mean $S=0.46$) from clean samples (mean $S=0.007$) by more than an order of magnitude, matching the synthetic separation. The gradient signature of shortcut behavior is therefore a property of training dynamics, not merely an artifact of small synthetic models.

\textbf{LLM-scale diagnosis.} The pilot identifies two structural constraints of autoregressive LLMs that are absent in the controlled synthetic regime. First, answer tokens carry both shortcut signal and the generation-supervision signal, so suppressing them at $\rho>0$ destabilizes training (perplexity rises five orders of magnitude); the appropriate scope of $\mathbf{g}_s^{\mathrm{ans}}$ is therefore the shortcut-relevant subset (e.g., numeric answer tokens on GSM8K) rather than all answer tokens. Second, batch-level gradient aggregation masks per-sample conflicts and reduces the rate at which correction fires, indicating that per-sample gating is required at this scale (Appendix~\ref{sec:appendix_llm_transfer}).

\textbf{Interpretation.} The detection geometry transfers cleanly across scale, while the correction step at LLM scale operates under autoregressive-generation constraints that the synthetic regime does not surface. Shortcut-aware training at LLM scale is therefore a constrained variant of the same gradient-alignment principle, with token-level supervision structure as an additional design axis. SART runs at $\sim 2.5\times$ the cost of plain SFT, with $\mathbf{g}_{\mathcal{V}}$ refreshed every $k=5$ steps (once per epoch in the LLM pilot).

\subsection{Summary}

The experiments converge on a consistent pattern: shortcut behavior is invisible at the loss level but clearly structured in gradient space; detection alone is insufficient because directional correction is what realizes generalization gains; the gradient signatures underlying \textsc{ShortcutScore} generalize across model scales; and scaling the correction layer to pretrained autoregressive LLMs requires per-sample conflict gating and numeric-only answer-token suppression. Together, these results support our formulation of shortcut-aware reasoning training as a problem of gradient alignment and directional correction.
